\section{The exact weighted theta operator and its domains}
\label{wt:sec:theta}

Throughout this section $t\in\C$ and
\[
 \lambda(x)=\log x,\qquad
 (U_tF)(x)=e^{t\lambda(x)^2/4}F(x),\qquad
 D=-x\frac{d}{dx}\quad(x>0).
\]
The factors $1/4$ here and $1/2$ in the squared norms below are
retained.  Let $V$ be the original space of even Schwartz functions
$\phi$ on $\R$ satisfying $\phi(0)=\int_\R\phi=0$, and let
\[
 \B=\left\{F\in C^\infty(0,\infty):
  p_{b,j}(F):=\sup_{x>0}x^b|D^jF(x)|<\infty
  \text{ for all }b\in\mathbb Z,\ j\geq0\right\}.
\]
The topology of $V$ is its original Schwartz subspace topology;
that of $\B$ is generated by the displayed seminorms.  Restriction
of an even function to $(0,\infty)$ is injective, so it causes no
ambiguity in the following definitions:
\begin{equation}\label{wt:eq:spaces}
 V_t=\{U_t(\phi|_{(0,\infty)}):\phi\in V\},\qquad
 \B_t=U_t\B.
\end{equation}
Both spaces have the transported topology.  In particular the
seminorms of $\B_t$ are exactly
$p^{(t)}_{b,j}(F)=p_{b,j}(U_{-t}F)$.  The source $V_t$ is a space
of positive-axis functions; it is not declared to be the original
even Schwartz source.

\begin{theorem}[Literal theta conjugacy]\label{wt:thm:theta}
The map $\Theta_t:V_t\longrightarrow\B_t$ defined by
$\Theta_t=U_t\Theta U_{-t}$ has the exact series
\begin{equation}\label{wt:eq:theta-series}
 (\Theta_tf)(x)
 =2\sum_{n=1}^{\infty}
 \exp\!\left(-\frac t2\log x\log n-\frac t4(\log n)^2\right)f(nx).
\end{equation}
The series and every differentiated series converge locally
uniformly on $(0,\infty)$.  With domains included, the Euler
operator and inversion become
\begin{equation}\label{wt:eq:Euler-reflection}
 D_t:=U_tDU_{-t}=D+\frac t2\log x:\B_t\longrightarrow\B_t,
 \qquad
 \mathcal I F(x):=x^{-1}F(x^{-1}),\qquad
 U_t\mathcal I=\mathcal I U_t.
\end{equation}
In particular $D_t^r=U_tD^rU_{-t}$ on $\B_t$ for every
integer $r\geq0$, where $D_t^r$ means operator composition.
\end{theorem}

\begin{proof}
Write $f=U_t\phi$ with $\phi\in V$ restricted to the positive
axis.  The $n$th term on the right side of
\eqref{wt:eq:theta-series}, including its exponential factor but
excluding the common factor $2$, is
\[
 e^{t(\log x)^2/4}\phi(nx),
\]
because
$(\log(nx))^2=(\log x)^2+2\log x\log n+(\log n)^2$.
On a compact interval $[a,b]\subset(0,\infty)$, each Euler
derivative of $\phi(nx)$ is $(D^j\phi)(nx)$.  Schwartz decay
bounds its absolute value by $C_{N,j,a}n^{-N}$ for any $N>1$.
Every derivative of $e^{t(\log x)^2/4}$ is bounded on $[a,b]$;
the product rule therefore gives summable majorants for every
differentiated series.  Ordinary derivatives are finite linear
combinations of Euler derivatives with smooth coefficients on
$[a,b]$, so the same assertion holds for them.  This proves both
the series identity and its smooth meaning.  The original map
$\Theta:V\to\B$ is continuous; hence its conjugate is continuous
between the transported spaces.

Direct differentiation gives
$D(U_{-t}F)=U_{-t}(DF+(t/2)(\log x)F)$, proving the first
operator identity, including its domain.  Induction proves the
power identity.  No commutative binomial formula is being used:
for example
\[
 D_t^2=D^2+t(\log x)D-\frac t2+\frac{t^2}{4}(\log x)^2,
\]
since $D(\log x)- (\log x)D=-I$.  Finally
$(\log(x^{-1}))^2=(\log x)^2$ proves the inversion identity
pointwise.  It transfers the original continuous involution
$\mathcal I$ of $\B$ to one of $\B_t$.
\end{proof}

Here and below the original closed-range theorem is the theorem
labelled \texttt{thm:closed} in
\texttt{ORIGINAL\_THETA\_CLOSED\_RANGE.tex}, whose definitions
of $V,\B,\Theta$ are the ones displayed above.  We use that proved
strong-space theorem, not a claim of closedness in an $L^2$ norm.

\begin{corollary}[The full strong quotient is transported]
\label{wt:cor:quotient}
The map $\Theta_t:V_t\to\B_t$ is a topological isomorphism
onto its closed range, and
\begin{equation}\label{wt:eq:strong-quotient}
 \overline U_t:\B/\Theta V\longrightarrow\B_t/\Theta_tV_t,
 \qquad [F]\longmapsto[U_tF]
\end{equation}
is a topological linear isomorphism, with inverse
$[G]\mapsto[U_{-t}G]$.  The bars here denote the induced quotient
map and do not denote closure.
\end{corollary}

\begin{proof}
By the definitions and Theorem~\ref{wt:thm:theta},
$\Theta_tV_t=U_t\Theta V$.  The map $U_t:\B\to\B_t$ is a
homeomorphism with inverse $U_{-t}$, so it carries the proved
closed subspace $\Theta V$ to a closed subspace.  The inverse on
the range is $U_t\Theta^{-1}U_{-t}$ and is continuous in the
transported source topology.  If $F-F'\in\Theta V$, then
$U_tF-U_tF'\in\Theta_tV_t$; the converse follows by applying
$U_{-t}$.  This proves that both displayed quotient maps are
well-defined and inverse.  Their continuity follows from the
definition of quotient topology: composition with the relevant
quotient projection is the composition of a continuous quotient
projection and $U_t$ or $U_{-t}$.
\end{proof}

\begin{proposition}[Transported Fourier map and both theta legs]
\label{wt:prop:Fourier-chain}
Let $\mathcal F$ be the original Fourier transform with kernel
$e^{-2\pi ixy}$ on the even Schwartz source $V$, and retain
$J=\mathcal I$, namely $JF(x)=x^{-1}F(x^{-1})$.  For every
$t\in\C$ the continuous involution on the transported source is
\begin{equation}\label{wt:eq:Fourier-transport}
 \mathcal F_t=U_t\mathcal F U_{-t}:V_t\longrightarrow V_t,
 \qquad \mathcal F_t^2=I.
\end{equation}
Its positive-axis integral, on its exact source domain, is
\begin{equation}\label{wt:eq:Fourier-kernel}
 (\mathcal F_tf)(x)=2e^{t(\log x)^2/4}
 \int_0^\infty e^{-t(\log y)^2/4}f(y)\cos(2\pi xy)\,dy.
\end{equation}
The identities with their source and target spaces are
\begin{equation}\label{wt:eq:Fourier-identities}
 \begin{gathered}
 J\Theta_t=\Theta_t\mathcal F_t:V_t\longrightarrow\B_t,
 \qquad \Theta_tD_t=D_t\Theta_t:V_t\longrightarrow\B_t,\\
 \mathcal F_tD_t=(1-D_t)\mathcal F_t:V_t\longrightarrow V_t,
 \qquad D_tJ=J(1-D_t):\B_t\longrightarrow\B_t.
 \end{gathered}
\end{equation}
In particular the two-leg differential and its source generator are
\begin{equation}\label{wt:eq:two-leg}
 \begin{gathered}
 \delta_t:V_t\oplus V_t\longrightarrow\B_t,
 \qquad\delta_t(f,g)=\Theta_tf-J\Theta_tg,\\
 \mathbb D_t=D_t\oplus(1-D_t),
 \qquad\delta_t\mathbb D_t=D_t\delta_t.
 \end{gathered}
\end{equation}
They are the exact transport of the original two-leg complex:
\begin{equation}\label{wt:eq:two-leg-conjugacy}
 \delta_t=U_t\delta_0(U_{-t}\oplus U_{-t}),\qquad
 \mathbb D_t=(U_t\oplus U_t)\mathbb D_0
                         (U_{-t}\oplus U_{-t}).
\end{equation}
The full kernel, range, and cokernel are
\begin{equation}\label{wt:eq:two-leg-kernel}
 \ker\delta_t=\{(\mathcal F_tg,g):g\in V_t\},\qquad
 \operatorname{im}\delta_t=\Theta_tV_t,\qquad
 \operatorname{coker}\delta_t=\B_t/\Theta_tV_t.
\end{equation}
The kernel identification with $V_t$ is a topological isomorphism
and carries its generator to $1-D_t$; the target quotient has
the descended generator $D_t$.
\end{proposition}

\begin{proof}
The original Fourier identities on $V$ are
$\mathcal F^2=I$, $J\Theta=\Theta\mathcal F$, and
$\mathcal FD=(1-D)\mathcal F$.  The first is Fourier inversion
on even functions, the second is the original Poisson identity
with both endpoint terms zero, and the third follows by integration
by parts in the Fourier integral.  They use exactly the original
Fourier kernel above.  The Fourier transform preserves $V$:
it preserves evenness, its value at zero is $\int_\R\phi=0$,
and its integral is $\phi(0)=0$ by inversion.  Its usual Schwartz
seminorm bounds give continuity on this source.  Thus conjugation
by $U_t$ proves \eqref{wt:eq:Fourier-transport} with the
transported topology.  Splitting the integral of an even function
into its two half axes gives the cosine kernel and the factor
$2$ in \eqref{wt:eq:Fourier-kernel}.  Here $U_{-t}f$ is exactly
the restriction of a member of $V$, so the integral converges
absolutely; no assertion about arbitrary positive-axis functions
is needed.

Conjugate the three original identities by $U_t$ and use
$JU_t=U_tJ$ to obtain the first and third identities in
\eqref{wt:eq:Fourier-identities}.  The original termwise
Euler identity $\Theta D=D\Theta$ and $DJ=J(1-D)$ give the
other two in exactly the same way.  All operators act on the
displayed spaces: in particular $D_t(V_t)\subset V_t$ follows
from $D(V)\subset V$ by conjugation.

For $(f,g)\in V_t\oplus V_t$, the identities just proved give
\[
 D_t\delta_t(f,g)
 =\Theta_tD_tf-J\Theta_t(1-D_t)g
 =\delta_t\mathbb D_t(f,g).
\]
The same conjugations give \eqref{wt:eq:two-leg-conjugacy}.
More explicitly, the continuous invertible source change
\[
 \mathcal C_t(f,g)=(f-\mathcal F_tg,g),\qquad
 \mathcal C_t^{-1}(a,g)=(a+\mathcal F_tg,g)
\]
satisfies
\[
 \delta_t=(\Theta_t,0)\mathcal C_t,\qquad
 \mathcal C_t\mathbb D_t=\mathbb D_t\mathcal C_t,
\]
because $\mathcal F_t(1-D_t)=D_t\mathcal F_t$ follows from
the third identity of \eqref{wt:eq:Fourier-identities}.
The map $\Theta_t$ is injective by
Corollary~\ref{wt:cor:quotient}.  Therefore
$\delta_t(f,g)=\Theta_t(f-\mathcal F_tg)$ vanishes exactly
when $f=\mathcal F_tg$.  Its image is $\Theta_tV_t$, since
$(f,0)$ attains every element of that image.  These prove all
of \eqref{wt:eq:two-leg-kernel}.  The graph map
$g\mapsto(\mathcal F_tg,g)$ and its inverse, projection to the
second coordinate, are continuous; applying $\mathbb D_t$
to this graph gives
$(\mathcal F_t(1-D_t)g,(1-D_t)g)$.  This proves its stated
generator.  The target generator descends because the image
is invariant, as follows from the second identity in
\eqref{wt:eq:Fourier-identities}.
\end{proof}

The transported spaces in Corollary~\ref{wt:cor:quotient} cannot silently be
replaced by the original ones.  The following examples use real
$t$.  Writing $u=\log x$,
the smooth function
$F(x)=\exp(-(1+u^2)^{3/4})$ belongs to $\B$: each Euler
derivative is bounded by a polynomial in $1+|u|$ times
$e^{-(1+u^2)^{3/4}}$, whose product with every $e^{b u}$
remains bounded on the real $u$ axis.
For $t>0$, $U_tF$ tends to infinity at both logarithmic ends,
so $U_tF\notin\B$.  For $t<0$, multiplication by $U_t$ maps
$\B$ into $\B$, by the product rule and Gaussian decay in $u$,
but the same $F$ is not in $U_t\B$, since $U_{-t}F\notin\B$.
Also $V_t$ need not equal the positive-axis restriction of $V$.
For $t\ne0$, select disjoint compact intervals in $(0,\infty)$
on which the continuous function $e^{t(\log x)^2/4}$ has disjoint
value ranges.  Nonnegative smooth bumps of integral one in those
intervals have difference of integral zero but a nonzero weighted
integral.  Its even extension lies in $V$, whereas the even
extension of its $U_t$ multiple has nonzero integral and does not.
These are explicit domain distinctions; the exact morphism
relating the domains and quotients is
\eqref{wt:eq:spaces}--\eqref{wt:eq:strong-quotient}.

\section{The moving Hilbert space and every tensor weight}
\label{wt:sec:Hilbert}

From this section onward, the Hilbert norms and their metric
derivatives use real $t$.  We use the column-coordinate inner-product convention
$\langle F,G\rangle=\int\overline F G$, conjugate-linear in
the first argument.  Define
\begin{equation}\label{wt:eq:Hilbert}
 \mathcal H_0=L^2((0,\infty),dx),\qquad
 \mathcal H_t=L^2((0,\infty),e^{-t(\log x)^2/2}dx).
\end{equation}
Multiplication gives a surjective isometry
\begin{equation}\label{wt:eq:isometry}
 U_t:\mathcal H_0\longrightarrow \mathcal H_t,\qquad
 \langle U_tF,U_tG\rangle_{\mathcal H_t}=\langle F,G\rangle_{\mathcal H_0},
 \qquad U_t^{-1}=U_{-t}:\mathcal H_t\longrightarrow \mathcal H_0.
\end{equation}
These identities follow by multiplication of the densities,
and prove surjectivity as well: for $G\in \mathcal H_t$, $U_{-t}G\in \mathcal H_0$
with exactly the same norm.  The original inclusion $\B\to \mathcal H_0$
is continuous, since
$\|F\|_{\mathcal H_0}^2\leq p_{0,0}(F)^2+p_{1,0}(F)^2$ by splitting
the integral at $1$.  Thus $\B_t\to \mathcal H_t$ is continuous too.
On the fixed space $\mathcal H_0$, multiplication by $U_t$ is unbounded
when $t>0$: unit vectors supported where $(\log x)^2\geq N$
have image norm at least $e^{tN/4}$.  Such compactly supported
vectors are in its multiplication domain.  For $t<0$ it is a
bounded contraction, but its inverse on its range is unbounded.
Neither assertion changes the isometry between the different
spaces in \eqref{wt:eq:isometry}.

For an integer $k\geq1$, write
\[
 d\mathbf x=dx_1\cdots dx_k,\qquad
 \Lambda_k(\mathbf x)=\sum_{j=1}^k(\log x_j)^2,\qquad
 \mathsf U_t=U_t^{\otimes k}.
\]
The Hilbert tensor products, identified with their product-measure
spaces, and the full transport are
\begin{equation}\label{wt:eq:tensor}
 \begin{split}
 \mathcal H_t^{\otimes k}
 &=L^2((0,\infty)^k,e^{-t\Lambda_k/2}d\mathbf x),\\
 (\mathsf U_tF)(\mathbf x)&=e^{t\Lambda_k(\mathbf x)/4}F(\mathbf x),\\
 \langle\mathsf U_tF,\mathsf U_tG\rangle_{\mathcal H_t^{\otimes k}}
 &=\langle F,G\rangle_{\mathcal H_0^{\otimes k}}.
 \end{split}
\end{equation}
The assertions first follow for products by iterated integration.
Finite sums of products are dense in the product $L^2$ spaces
(rectangular simple functions are dense), so isometry extends
them to the completed tensor products.  Its inverse is again
multiplication by $e^{-t\Lambda_k/4}$.
The fixed-$\mathcal H_0^{\otimes k}$ pairing instead is
\begin{equation}\label{wt:eq:fixed-form}
 A_t(F,G):=\langle\mathsf U_tF,\mathsf U_tG\rangle_{\mathcal H_0^{\otimes k}}
 =\int_{(0,\infty)^k}e^{t\Lambda_k/2}\overline{F(\mathbf x)}
 G(\mathbf x)\,d\mathbf x,
\end{equation}
on the functions for which it is finite.  In particular the
weight is the sum of the $k$ logarithmic squares, not the square
of the sum of the logarithms.
On smooth transported tensor sources one also has
\[
 \mathsf U_t\Bigl(\sum_{j=1}^kD_j\Bigr)\mathsf U_{-t}
 =\sum_{j=1}^kD_j+\frac t2\sum_{j=1}^k\log x_j.
\]
The tensor theta series retains a separate integer in each
coordinate: its coefficient on $f(n_1x_1,\ldots,n_kx_k)$ is
\[
 2^k\exp\!\left(-\frac t2\sum_{j=1}^k\log x_j\log n_j
                 -\frac t4\sum_{j=1}^k(\log n_j)^2\right).
\]
For finite sums of transported product sources, convergence and
the tensor conjugacy follow coordinate by coordinate from
Theorem~\ref{wt:thm:theta}; the product of the summable majorants
justifies every iterated sum.

\section{A differentiable attained metric on the actual finite source}
\label{wt:sec:metric}

The double-Gaussian Euler class used here is the class $\Epi$
of smooth $F$ for which
\begin{equation}\label{wt:eq:class}
 \sup_{x>0}e^{a(x^2+x^{-2})}|D^jF(x)|<\infty
 \quad(j\geq0,\ 0<a<\pi).
\end{equation}
Let $\Epi^{\otimes_{\mathrm{alg}} k}$ denote finite sums of
products of $k$ such functions, not a completed tensor space.

\begin{lemma}[Domination on the original source]
\label{wt:lem:domination}
Multiplication by $U_t$ preserves $\Epi$ for every real $t$.
For $F,G\in\Epi^{\otimes_{\mathrm{alg}} k}$ the integral
in \eqref{wt:eq:fixed-form} is finite for every real $t$, is
real analytic in $t$, and for every integer $r\geq0$ satisfies
\begin{equation}\label{wt:eq:form-derivative}
 \frac{d^r}{dt^r}A_t(F,G)
 =2^{-r}\int\Lambda_k^r e^{t\Lambda_k/2}\overline F G\,d\mathbf x.
\end{equation}
The same statements hold simultaneously on any fixed finite
dimensional subspace of this class.
\end{lemma}

\begin{proof}
For $0<a<a'<\pi$, the function
\[
 (1+|\log x|)^N
 \exp\!\left(T(\log x)^2-(a'-a)(x^2+x^{-2})\right)
\]
is bounded on $(0,\infty)$ for every fixed $N\geq0$ and $T\geq0$.
Indeed in $u=\log x$ the negative term is
$-(a'-a)(e^{2u}+e^{-2u})$, which dominates every polynomial
in $u$ as $u\to\pm\infty$; boundedness on a compact interval
follows from continuity.  Each Euler derivative of $U_tF$ is
$U_t$ times a finite sum of Euler derivatives of $F$ multiplied
by polynomials in $\log x$, by the product rule.  Applying the
bound with $a'$ proves \eqref{wt:eq:class} for $U_tF$, uniformly
on compact $t$ intervals.

A product source, and hence any finite sum of product sources,
is bounded in absolute value by
$C\exp(-a\sum_j(x_j^2+x_j^{-2}))$ for any fixed $0<a<\pi$.
For complex $z$ in a compact disk the derivatives of
$e^{z\Lambda_k/2}\overline F G$ of every prescribed finite
order are therefore dominated by a constant times
$\exp(-a\sum_j(x_j^2+x_j^{-2}))$, using the preceding bound
coordinate by coordinate.  This is integrable: near infinity
it is bounded by a Gaussian and near zero by $e^{-a/x^2}$,
which is bounded on an interval of finite length.  Differentiating
under the integral is justified by this majorant.  It also proves
complex holomorphy for the integral with $z$ in place of $t$,
by dominated difference quotients, hence real analyticity and
\eqref{wt:eq:form-derivative}.  A basis of a finite dimensional
subspace requires only finitely many such bounds.
\end{proof}

\begin{theorem}[Exact quotient metric and its derivative]
\label{wt:thm:metric}
Let $W\subset\Epi^{\otimes_{\mathrm{alg}} k}$ be finite
dimensional, and let $\mathcal J:W\to Y$ be a surjective linear
map to a finite dimensional space.  Keep this map fixed on $W$.
On $W_t=\mathsf U_tW$ the transported map is
\[
 \mathcal J_t=\mathcal J\mathsf U_{-t}:W_t\longrightarrow Y,
 \qquad \ker\mathcal J_t=\mathsf U_t(\ker\mathcal J).
\]
Give $W_t$ either its induced moving-$\mathcal H_t^{\otimes k}$ norm or
its induced fixed-$\mathcal H_0^{\otimes k}$ norm.  The former induces
on $Y$ the original, constant quotient metric.  For the latter,
choose a basis $f_0,\ldots,f_{d-1}$ of $W$ and column coordinates
on $Y$.  Let $E$ be the surjective matrix of $\mathcal J$, and put
\begin{equation}\label{wt:eq:matrix-formula}
 \begin{split}
 (K_t)_{nm}&=A_t(f_n,f_m),\\
 G_t&=(E K_t^{-1}E^*)^{-1},\\
 R_t&=K_t^{-1}E^*G_t,\qquad ER_t=I.
 \end{split}
\end{equation}
Then the fixed-space attained quotient norm of $a\in Y$ is
$a^*G_ta$, and its unique least-norm lift in $W_t$ is
\begin{equation}\label{wt:eq:lift}
 \mathsf U_tw_{t,a},\qquad
 w_{t,a}=\sum_{n=0}^{d-1}(R_ta)_nf_n.
\end{equation}
All these matrices are real analytic for real $t$, and
\begin{equation}\label{wt:eq:derivative}
 G'_t=R_t^*K'_tR_t,\qquad
 a^*G'_tb
 =\frac12\int\Lambda_ke^{t\Lambda_k/2}
       \overline{w_{t,a}}w_{t,b}\,d\mathbf x.
\end{equation}
In particular $G'_t$ is positive definite whenever $Y\ne\{0\}$.
At $t=0$ the functions $w_{0,a}$ in this formula are precisely
the original least-norm lifts; their first derivatives are not
needed to evaluate the metric derivative.
\end{theorem}

\begin{proof}
The moving-space assertion follows from the bijection
$w\mapsto\mathsf U_tw$ between the two fibers with value $a$
and the exact isometry \eqref{wt:eq:tensor}.
For the fixed-space assertion, $K_t$ is positive definite:
its quadratic form on a nonzero coefficient vector is the
positive-density integral of the squared absolute value of a
nonzero source function.  Linear independence makes that function
nonzero in $L^2$; the functions are smooth, so a nonzero function
is nonzero on an open set.  For $a\ne0$, $E^*a\ne0$ because
$E$ is onto, and therefore
$a^*EK_t^{-1}E^*a>0$.  Both inverses in
\eqref{wt:eq:matrix-formula} thus exist.

One has $ER_t=I$ and $K_tR_t=E^*G_t$.  Every coefficient lift
of $a$ is uniquely $R_ta+z$ with $Ez=0$.  Its squared norm is
\[
 (R_ta+z)^*K_t(R_ta+z)=a^*G_ta+z^*K_tz,
\]
since $R_t^*K_tz=G_tEz=0$ and
$R_t^*K_tR_t=G_t$.  This proves uniqueness, the metric, and
\eqref{wt:eq:lift}.  Lemma~\ref{wt:lem:domination} gives real
analyticity of $K_t$; invertibility at every real $t$ and the
finite matrix inverse formula give local real analyticity of
$G_t$ and $R_t$ there.

Differentiating $K_tK_t^{-1}=I$ yields
$(K_t^{-1})'=-K_t^{-1}K'_tK_t^{-1}$.  Differentiating the
inverse defining $G_t$ then gives
\[
 G'_t=G_tEK_t^{-1}K'_tK_t^{-1}E^*G_t=R_t^*K'_tR_t.
\]
The integral follows from \eqref{wt:eq:form-derivative}.
Alternatively, differentiating $R_t^*K_tR_t=G_t$ gives the
same result because $ER'_t=0$ and
$(R'_t)^*K_tR_t=(ER'_t)^*G_t=0$, together with its adjoint.
Finally $\Lambda_k$ vanishes only at the single point
$(1,\ldots,1)$, a null set.  If $a\ne0$, then
$\mathcal Jw_{t,a}=a$ implies $w_{t,a}\ne0$; its weighted
logarithmic-square integral is strictly positive.  This proves
the last assertion and also strict increase of the quotient
metric with $t$ on every nonzero vector.
\end{proof}

For comparison with the convention linear in the first argument,
the last formula at $k=1,t=0$ reads exactly
\begin{equation}\label{wt:eq:linear-first}
 \left.\frac d{dt}\right|_{t=0}q_t^{\mathrm{lin}}(a,b)
 =\frac12\int_0^\infty(\log x)^2w_{0,a}(x)
                         \overline{w_{0,b}(x)}\,dx,
 \qquad q_t^{\mathrm{lin}}(a,b)=b^*G_ta.
\end{equation}
Thus changing the inner-product convention transposes the indexed
Gram matrix; it does not change the positive sign or the factor
$1/2$.

\begin{proposition}[Schur complement in original minimal lifts]
\label{wt:prop:Schur}
Choose a basis $a_1,\ldots,a_q$ of $Y$, let $u_i=w_{0,a_i}$
be its original minimal lifts, and choose a basis
$k_1,\ldots,k_{d-q}$ of $\ker\mathcal J$.  In the basis
$(u_1,\ldots,u_q,k_1,\ldots,k_{d-q})$, write the column-convention
Gram matrix of $A_t$ as
\[
 \begin{pmatrix}B_t&C_t\\C_t^*&\mathsf H_t^{\mathrm{ker}}\end{pmatrix}.
\]
For a nonzero kernel the quotient Gram and lift are
\begin{equation}\label{wt:eq:Schur}
 G_t^{\mathrm{Schur}}=B_t-C_t(\mathsf H_t^{\mathrm{ker}})^{-1}C_t^*,
 \qquad
 s_t(a_i)=u_i-\sum_\alpha
   ((\mathsf H_t^{\mathrm{ker}})^{-1}C_t^*)_{\alpha i}k_\alpha.
\end{equation}
For zero kernel they are respectively $B_t$ and $u_i$.
Original minimality gives $C_0=0$, and consequently
\[
 (G_0^{\mathrm{Schur}})'_{ij}
 =\frac12\int\Lambda_k\overline{u_i}u_j\,d\mathbf x.
\]
\end{proposition}

\begin{proof}
For a coefficient column $a$ in the chosen quotient basis the
fiber has coefficient columns $(a,z)$.  Its squared norm equals
\[
 a^*(B_t-C_t(\mathsf H_t^{\mathrm{ker}})^{-1}C_t^*)a
 +\bigl(z+(\mathsf H_t^{\mathrm{ker}})^{-1}C_t^*a\bigr)^*
 \mathsf H_t^{\mathrm{ker}}
 \bigl(z+(\mathsf H_t^{\mathrm{ker}})^{-1}C_t^*a\bigr).
\]
Positive definiteness of the kernel block proves the two
formulas.  Since each $u_i$ is orthogonal to the kernel at zero,
$C_0=0$.  Every term in the derivative of the correction
$C_t(\mathsf H_t^{\mathrm{ker}})^{-1}C_t^*$ therefore contains a zero
factor at zero, leaving $B'_0$.  Lemma~\ref{wt:lem:domination}
evaluates this derivative as displayed.
\end{proof}

\section{Specialization to the complete arithmetic packet}
\label{wt:sec:packet}

Retain the original complete-order packet and source
\[
 h(s)=\prod_{\rho\in Z}(s-\rho)^{m_\rho},\qquad
 g(s)=2\xi(s),\qquad \Mell F_h=g/h,
 \qquad \upsilon=j_h(g/h)\in E_h=\C[s]/(h).
\]
Here $j_h$ uses the Taylor coefficients
$f^{(r)}(\rho)/r!$ for $0\leq r<m_\rho$ in each local factor.
The source theorem \texttt{thm:divisor} in the present derivation
proves $F_h\in\Epi$; the jet theorem \texttt{thm:jets} proves
that $\upsilon$ is a unit.  These are proved properties of the
actual source, not replacements by a selected comparison source.
Set, exactly as in SA1--SA5 and SA37 of
\texttt{NATIVE\_SPECTRAL\_ACTION.tex},
\begin{gather*}
 w_h(v)=\frac{|(g/h)(1/2+iv)|^2}{2\pi},\qquad
 m=w_h^{*k},\qquad c=k/2,\\
 \chi(S)=\prod_\kappa(S-\kappa)^{e_\kappa},\qquad
 e_\kappa=\max_{\rho_1+\cdots+\rho_k=\kappa}
      \left(1+\sum_{j=1}^k(m_{\rho_j}-1)\right),\\
 C=\C[S]/(\chi),\qquad q=\deg\chi,\qquad L\geq q-1.
\end{gather*}
Every sum collision and original multiplicity is retained in
$\chi$.

\begin{lemma}[The original measure and its orthogonal polynomials]
\label{wt:lem:measure}
The function $w_h$ is real, nonnegative, even, positive almost
everywhere, and has every absolute polynomial moment.  Its
mass is finite and strictly positive.  The convolution $m=w_h^{*k}$
has the same properties; when $k\geq2$ it is positive at every
real argument.  For each integer $n\geq0$ there exists a unique
real monic polynomial $Q_n$ of degree $n$ orthogonal to every
polynomial of smaller degree for $m(u)\,du$.  It satisfies
$Q_n(-u)=(-1)^nQ_n(u)$ and has a finite strictly positive
squared norm $\omega_n$.
\end{lemma}

\begin{proof}
The function $R_h=g/h=\Mell F_h$ is a nonzero entire function.
Nonzero follows, for example, from $hR_h=g$ and $g\not\equiv0$.
An entire function not identically zero has isolated zeros:
at any zero take the first nonzero coefficient in its Taylor
series and factor out its finite order.  Thus its zeros on
$s=1/2+iv$ form a set with finitely many points on each compact
interval of $v$, and hence a countable null set.  This proves
$w_h(v)>0$ almost everywhere.

Since $f_0$ is real-valued, its entire Mellin representation
gives $g(\overline s)=\overline{g(s)}$.  The packet is stable
under conjugation with the actual zero orders, so its monic
polynomial satisfies $h(\overline s)=\overline{h(s)}$.
Division away from its finitely many roots, followed by
continuation across them, gives
$R_h(\overline s)=\overline{R_h(s)}$.  Consequently
$|R_h(1/2-iv)|=|R_h(1/2+iv)|$, proving evenness of $w_h$.

For completeness, put $y=\log x$.  The function
$e^{y/2}F_h(e^y)$ is Schwartz in $y$, as is each of its
derivatives: the Euler bounds for $F_h$ and the product rule
bound them by fixed powers of $e^y,e^{-y},|y|$ times
$e^{-a(e^{2y}+e^{-2y})}$ for any fixed $0<a<\pi$.
The same holds after multiplication by any power of $y$.
The integral
\[
 R_h(1/2+iv)=\int_\R e^{y/2}F_h(e^y)e^{ivy}\,dy
\]
can therefore be differentiated in $v$ and integrated by parts
in $y$ arbitrarily many times, with zero boundary terms.
Each resulting integrand is integrable.  This proves that
$R_h(1/2+iv)$ and all its $v$ derivatives decrease faster
than every inverse power of $1+|v|$.  By the product rule,
$w_h(v)=R_h(1/2+iv)\overline{R_h(1/2+iv)}/(2\pi)$ is
Schwartz as well.  In particular it is bounded and all numbers
\[
 M_N=\int_\R |v|^N w_h(v)\,dv\quad(N\geq0)
\]
are finite, with $M_0>0$ by the already proved positivity.

Each iterated convolution is finite at every argument, since
$w_h$ is bounded and integrable; induction gives the bound
$\|w_h^{*k}\|_\infty\leq\|w_h\|_\infty M_0^{k-1}$.
For fixed $u$, the integrand $w_h(v)w_h(u-v)$ is positive
for almost every $v$: the two exceptional sets are null,
including the translated reflected zero set.  Its integral is
therefore strictly positive.  The same argument applied to
$w_h(v)w_h^{*(k-1)}(u-v)$ proves positivity at every $u$ for
all $k\geq2$.  Evenness follows at each stage by replacing
$v$ with $-v$ in the convolution.  Tonelli's theorem and the
definition of convolution give
\[
 \int_\R m(u)\,du=M_0^k,\qquad
 \int_\R |u|^N m(u)\,du
 =\int_{\R^k}\left|\sum_{j=1}^kv_j\right|^N
                         \prod_{j=1}^kw_h(v_j)\,d\mathbf v.
\]
For each integer $N\geq1$, convexity gives
$|\sum_jv_j|^N\leq k^{N-1}\sum_j|v_j|^N$; the last
integral is consequently at most $k^NM_NM_0^{k-1}$.
Thus all required convolution moments are finite.

For any nonzero polynomial $P$, its real zero set is finite,
and $m>0$ almost everywhere.  Hence
$\int |P(u)|^2m(u)\,du>0$, while the moment bounds make it
finite.  The moment pairing is therefore positive definite
on each finite polynomial space.  Gram--Schmidt on
$1,u,u^2,\ldots$, retaining leading coefficient one, produces
the unique monic orthogonal polynomials; every denominator is
strictly positive by this pairing.  The moments are real,
so all resulting coefficients are real.  Finally
$(-1)^nQ_n(-u)$ is another monic degree-$n$ polynomial with
the same orthogonality by evenness of $m$.  Uniqueness proves
the asserted parity.  Positivity and finiteness of $\omega_n$
are the same positive-definite pairing applied to $Q_n$.
\end{proof}

Using these actual orthogonal polynomials, put
\[
 \omega_n=\int_\R Q_n(u)^2m(u)\,du,\qquad
 p_n(S)=i^nQ_n((S-c)/i),\qquad b_n=[p_n]_\chi.
\]
In particular $\omega_0=(\int_\R w_h)^k$ is the original mass;
it is not reset to one.  Define the actual polynomial source map
and its original orthonormal basis by
\begin{equation}\label{wt:eq:actual-source}
 \mathcal AP=P(D_1+\cdots+D_k)F_h^{\otimes k},\qquad
 P_n(S)=\frac{i^{-n}p_n(S)}{\sqrt{\omega_n}},\qquad
 f_n=\mathcal AP_n\quad(0\leq n\leq L),\qquad
 W_L=\mathcal A\C[S]_{\leq L}.
\end{equation}
Each $f_n$ belongs to $\Epi^{\otimes_{\mathrm{alg}} k}$,
because expanding its finite polynomial in $\sum_jD_j$ gives
a finite sum of products of Euler derivatives of $F_h$.

For clarity, the Hilbert identity making this the actual basis
is the original Mellin--Parseval identity with all factors:
\begin{equation}\label{wt:eq:parseval}
 \begin{split}
 \|\mathcal AP\|_{\mathcal H_0^{\otimes k}}^2
 &=\int_{\R^k}\left|P\left(c+i\sum_{j=1}^kv_j\right)\right|^2
              \prod_{j=1}^kw_h(v_j)\,d\mathbf v\\
 &=\int_\R|P(c+iu)|^2m(u)\,du.
 \end{split}
\end{equation}
Indeed integration by parts on the double-Gaussian source gives
$\Mell(D_jF)(\mathbf s)=s_j\Mell F(\mathbf s)$, with zero
boundary terms.  With $y=\log x$, the critical-line Mellin
integral is the Fourier integral of $e^{y/2}F(e^y)$, whose
$L^2(dy)$ norm is exactly $\|F\|_{L^2(dx)}$; the Fourier
Parseval factor is $1/(2\pi)$ for the frequency $v$ used here.
Iterating gives the first identity; the defining convolution
integral gives the second.  Since $P_n(c+iu)=Q_n(u)/\sqrt{\omega_n}$,
the $f_n$ are orthonormal.  Positivity of the original measure
also proves injectivity of $\mathcal A$ on polynomials: a nonzero
polynomial has finitely many real-line zeros, so its squared
norm in \eqref{wt:eq:parseval} is positive.

Define the fixed surjection
\begin{equation}\label{wt:eq:actual-jet}
 \mathcal J:W_L\longrightarrow C,\qquad
 \mathcal J(\mathcal AP)=[P]_\chi,\qquad
 E=\left[\frac{i^{-n}b_n}{\sqrt{\omega_n}}\right]_{n=0}^L.
\end{equation}
It is well-defined by injectivity of $\mathcal A$.  It is onto
because the first $q$ monic polynomials form a basis after
reduction.  The actual tensor jet is connected to this quotient
by the injective map
\begin{equation}\label{wt:eq:eta}
 \eta:C\longrightarrow E_h^{\otimes k},\qquad
 \eta[P]=\upsilon^{\otimes k}P(s_1+\cdots+s_k),\qquad
 j_h^{\otimes k}\Mell(\mathcal AP)=\eta\mathcal J(\mathcal AP).
\end{equation}
The last equality is the Euler--Mellin identity just proved,
followed by finite Taylor reduction.  To verify injectivity,
in a local tensor factor put $s_j=\rho_j+z_j$ with
$z_j^{m_{\rho_j}}=0$.  The largest nonzero power of $\sum_jz_j$
is $\sum_j(m_{\rho_j}-1)$: its coefficient on
$\prod_jz_j^{m_{\rho_j}-1}$ is the nonzero multinomial
coefficient
$(\sum_j(m_{\rho_j}-1))!/\prod_j(m_{\rho_j}-1)!$.
All higher powers vanish.  Thus the annihilator in that factor
is $(S-\sum_j\rho_j)^{1+\sum_j(m_{\rho_j}-1)}$.
Taking the least common multiple over factors yields exactly
$\chi$, with the displayed maximum at collisions.  Finally
multiplication by the unit $\upsilon^{\otimes k}$ preserves
the kernel, proving injectivity of $\eta$.

Apply Theorem~\ref{wt:thm:metric} to this exact $W_L$ and
$\mathcal J$.  In the original basis \eqref{wt:eq:actual-source},
the complete result is
\begin{equation}\label{wt:eq:actual-derivative}
 \begin{gathered}
 (K_t)_{nm}=\int e^{t\Lambda_k/2}\overline{f_n}f_m\,d\mathbf x,
 \qquad K_0=I,\qquad
 G_t=(EK_t^{-1}E^*)^{-1},\\
 G_0=G=(EE^*)^{-1},\qquad R_0=R=E^*G,\\
 G'_0=R^*K'_0R,\qquad
 (K'_0)_{nm}=\frac12\int\left(\sum_{j=1}^k(\log x_j)^2\right)
                          \overline{f_n}f_m\,d\mathbf x,\\
 a^*G'_0b=\frac12\int\left(\sum_{j=1}^k(\log x_j)^2\right)
 \overline{\left(\sum_n(Ra)_nf_n\right)}
             \left(\sum_m(Rb)_mf_m\right)\,d\mathbf x.
 \end{gathered}
\end{equation}
At any real $t$ use $R_t=K_t^{-1}E^*G_t$ and the full
$e^{t\Lambda_k/2}$ weight as in \eqref{wt:eq:derivative}.
The physical image $\eta(C)$ receives this metric by
$\langle\eta a,\eta b\rangle_t=a^*G_tb$.  Its transported
source map is $\eta\mathcal J\mathsf U_{-t}$ on
$\mathsf U_tW_L$, and its minimal physical lift is exactly
$\mathsf U_t\sum_n(R_ta)_nf_n$.  Thus the same actual jets,
the same collision algebra, the same original mass, and their
attained source lifts remain linked by explicit maps at every
real $t$.

\section{The precise finite-dimensional scope}
\label{wt:sec:scope}

The strong quotient \eqref{wt:eq:strong-quotient} and the finite
Hilbert quotient \eqref{wt:eq:actual-derivative} are both proved,
and their source maps are explicit.  Neither construction says
that evaluation of Mellin jets extends to a bounded map on all
of $\mathcal H_0$.  In fact it does not.  On $C_c^\infty(0,\infty)$
the $r$th normalized Mellin jet at $\rho\in\C$ is
\[
 \ell_{\rho,r}(F)=\frac1{r!}\int_0^\infty
       F(x)x^{\rho-1}(\log x)^r\,dx.
\]
If this functional were bounded for the $L^2(dx)$ norm, its
Riesz representative would coincide on every compact subinterval
with $x^{\overline\rho-1}(\log x)^r/r!$, by testing against
all compactly supported smooth functions there.  Hence that
function would belong to $L^2(0,\infty)$.  Its squared absolute
value, however, has integral
\[
 \frac1{(r!)^2}\int_0^\infty
   x^{2\operatorname{Re}\rho-2}|\log x|^{2r}\,dx=\infty.
\]
If $\operatorname{Re}\rho>1/2$ it diverges at infinity;
if $\operatorname{Re}\rho<1/2$ it diverges at zero; at equality
the substitution $u=\log x$ gives the divergent integral of
$|u|^{2r}$ on $\R$.  This contradiction proves unboundedness.
The compactly supported tests lie in $\Epi$, so the obstruction
already occurs on the concrete source class, when it is given
only the ambient $L^2$ norm.  In contrast every jet map on the
finite space $W_L$ is continuous, and its least-norm quotient
is exactly the finite attained metric proved above.  A bounded
extension of the transported jet to all of $\mathcal H_t$ would, by the
isometry $U_t$, imply the forbidden bounded original extension.
This establishes the scope by an exact isometry and an explicit
unbounded functional, not by assuming a missing global quotient.
